\scanpage{112}
% archive sheet number: 111
\[
\begin{aligned}
 D_0=H_1\cap H_2:\quad &y=z,\ x=-\alpha'y\quad\text{ou}\quad y=z=\alpha x\\
 &\text{c'est la droite homogène de \unsure{vect.-dir.}}\\
 &(1,\alpha,\alpha)\quad\text{ou}\quad(-\alpha',1,1)
\end{aligned}
\]
\[
 D_1=H_2\cap H_0\quad
 \left\{\begin{aligned}
 &\text{\unsure{2 Nilpotent|2 nilpotent}:}\quad D_1=\varnothing,\ (\text{i.e. }H_0\parallel H_2)\\
 &\deleted{\unsure{y}}\\
 &\text{2 inv.}\qquad\deleted{\unsure{D_1}}\quad y=z=\frac12-x
 \end{aligned}\right.
\]
\[
 D_2=H_0\cap H_1\qquad
 \left\{\begin{aligned}
 y&=\alpha+(1-3\alpha)x\\
 z&=(1-\alpha)-3(1-\alpha)x
 \end{aligned}\right.
\]

Se donner une réal. géom. régulière de $\Pi$ revient au $=$ que de se donner une représentation affine de rang 3 de $\Gamma$ \inserted{(\unsure{dont|dans} $E$ l'esp. affine)}, \uline{et} un \deleted{élément} \inserted{\unsure{point}} $s_0$ de $E$ \unsure{telle|tel} que
\[
 \sigma_1s_0=\sigma_2s_0=s_0\quad\text{i.e.}\quad s_0\in E^{\sigma_1}\cap E^{\sigma_2}
\]
et \unsure{telle|tel} que de \uline{plus} les quatre pts
\[
 s_0,\ s_1=\sigma_0s_0,\ s_2=\sigma_1s_1=\sigma_1\sigma_0s_0,\ s_3=\sigma_2s_2=\sigma_2\sigma_1\sigma_0s_0
\]
soient affinement indépendants. \deleted{\unsure{Tous les}} \deleted{points sont conjugués \unsure{sous}} \unsure{À} \uline{deux} tels $s_0,s'_0$, il y a un autom. unique de $E$ commutant à $\Gamma$ qui transforme $s_0$ en $s'_0$.

% margin note written sideways
% slanted
\margin{Pour une représentation affine fidèle \unsure{et|?} (i.e. non triviale sur l'antipodisme ni sur le groupe $\unsure{\mathfrak{A}_5|\mathfrak{A}}$) sur chaque fibre, \deleted{\unsure{les}} conditions d'existence de $s_0$ sont-elles vérifiées ? Cela liquiderait la question de la classification des repr. affines du groupe de l'icosaèdre\unsure{\ldots}}
\margin{?}
